\section{Discussion}
Feasibility is not deployability, and the gap is a \emph{dual obstruction}. The feasible, post-quantum-sound proof, the bare STARK receipt, is not zero-knowledge (Horn~1); and the zero-knowledge wrap the toolchain ships, though it now fits the $24\,\mathrm{GB}$ budget, is a Groth16/PLONK SNARK over a pairing-friendly curve and hence not post-quantum (Horn~2), retroactively de-anonymizing every stored transcript to a future quantum adversary. No proof the stack produces is at once zero-knowledge and post-quantum-sound: the obstruction is a primitive limit, not feasibility. For the anonymous-credential construction BDEC~\cite{10.1007/978-981-96-0957-4_3} that PLUM instantiates, we further prove that \emph{everlasting} anonymity cannot be established in the standard model for its hash-committed showings~\cite{cgh2004rom}, escapable only by a post-quantum statistically-hiding commitment to the witness. The positive results are conditional, resting on five named open assumptions (Griffin-AIR soundness, transcript simulatability, the quantum-random-oracle lift, lookup binding, and signature key-privacy).
